% 1.3 群表示
\section{群表示}

现在我们从 1.2 节的抽象群论转向本书真正关心的专题——群的表示。因为本书的主题是固体研究中出现的那些群的表示，而非抽象群论本身的发展。对纯粹数学家而言，更高深的抽象群论中确有不少有趣的课题（例如中心扩张与上同调群），它们能帮助我们理解所关心的群；但对面向应用数学家与理论物理读者的书来说，深入这些课题将是一段漫长而难以消化的弯路。本节继续给出若干初等的定义和定理，并尽可能仍以群 $\mathbf{G}_6^2$ 为例加以说明。

\textbf{定义 1.3.1.} \textit{矩阵群.} 一个\textit{矩阵群} $\boldsymbol{\Delta}$ 是由非奇异矩阵组成的群。若其全部矩阵都是幺正的，则称为\textit{幺正矩阵群}。

以下只涉及有限阶的矩阵群与有限维的矩阵。

\textbf{定义 1.3.2.} \textit{等价.} 若存在非奇异矩阵 $\mathbf{S}$ 使 $\mathbf{D}_{1} = \mathbf{S}\mathbf{D}_{2}\mathbf{S}^{-1}$，则称矩阵 $\mathbf{D}_{1}$ 与 $\mathbf{D}_{2}$ 共轭。

若存在非奇异矩阵 $\mathbf{S}$ 使 $\boldsymbol{\Delta}_{1} = \mathbf{S}\boldsymbol{\Delta}_{2}\mathbf{S}^{-1}$，则称矩阵群 $\boldsymbol{\Delta}_{1}$ 与 $\boldsymbol{\Delta}_{2}$ \textit{等价}。注意这不仅要求两个群的矩阵能够成对共轭，而且所有共轭都必须由\textit{同一个}矩阵 $\mathbf{S}$ 实现。这是很强的条件：由此立即推出 $\boldsymbol{\Delta}_{1}$ 与 $\boldsymbol{\Delta}_{2}$ 必同构。反之不然，因为 $\boldsymbol{\Delta}_{1}$ 与 $\boldsymbol{\Delta}_{2}$ 可以同构却维数不同，此时显然不等价。

\textbf{定理 1.3.1.} \textit{每个矩阵群都等价于某个幺正矩阵群。}

处理矩阵时将使用以下记号：$\mathbf{D}^{T}$ 表示 $\mathbf{D}$ 的转置；$\mathbf{D}^{\*}$ 表示 $\mathbf{D}$ 的复共轭；$\mathbf{D}^{\dagger}\,[=(\mathbf{D}^{\*})^{T}]$ 表示 $\mathbf{D}$ 的厄米共轭；$\tilde{\mathbf{D}}\,[=(\mathbf{D}^{-1})^{T}]$ 表示 $\mathbf{D}$ 的逆步矩阵（contragredient）；$\dim \mathbf{D}$ 表示 $\mathbf{D}$ 的维数。

\textbf{定义 1.3.3.} \textit{迹、特征标.} 矩阵 $\mathbf{D}$ 的\textit{迹}（trace）是其对角元之和，记作 $\mathrm{tr}\ \mathbf{D}$。

矩阵群 $\boldsymbol{\Delta}$ 的\textit{特征标}（character）是定义在所有元素 $\mathbf{D} \in \boldsymbol{\Delta}$ 上的函数 $\chi$：$\chi(\mathbf{D}) = \mathrm{tr}\ \mathbf{D}$。

\textbf{定理 1.3.2.} \textit{若 $\boldsymbol{\Delta}$ 是具有单位元 $\mathbf{E}$ 的矩阵群，则 $\chi(\mathbf{E}) = \dim \mathbf{E}$；并且，若 $\mathbf{C}$ 与 $\mathbf{D}$ 属于 $\boldsymbol{\Delta}$ 的同一共轭类，则 $\chi(\mathbf{C}) = \chi(\mathbf{D})$。}

\textbf{定理 1.3.3.} \textit{两个矩阵群同构且具有相同特征标，当且仅当它们等价。}

\textbf{定义 1.3.4.} \textit{群的表示.} 群 $\mathbf{G}$ 的一个\textit{表示}（representation）是 $\mathbf{G}$ 到某群 $\mathbf{T}$ 上的同态 $\gamma$，$\mathbf{T}$ 由作用在复数域上有限维向量空间 $V$ 中的非奇异线性算符组成。对一切 $G \in \mathbf{G}$ 记 $\gamma G = \mathbf{T}_{G}$。

由定义 1.3.4，当 $\gamma$ 是表示时即有：

\begin{quote}
(i) 对一切 $G_{1}, G_{2} \in \mathbf{G}$ 与一切 $\mathbf{x} \in V$，$\mathbf{T}_{G_{1}}(\mathbf{T}_{G_{2}}\mathbf{x}) = \mathbf{T}_{G_{1}G_{2}}\mathbf{x}$；
(ii) 对一切 $\mathbf{x} \in V$，$\mathbf{T}_{E}\mathbf{x} = \mathbf{x}$，即 $\mathbf{T}_{E}$ 是恒等算符；
(iii) 对一切 $G \in \mathbf{G}$ 与一切 $\mathbf{x} \in V$，$\mathbf{T}_{G}^{-1}\mathbf{x} = \mathbf{T}_{G^{-1}}\mathbf{x}$。
\end{quote}

若 $\gamma$ 还是同构，则称该表示为\textit{忠实的}（faithful）。

设我们取由 $d$ 个线性无关向量 $\mathbf{x}_{1}, \mathbf{x}_{2}, \ldots, \mathbf{x}_{d}$ 组成、张成空间 $V$ 的基 $\langle \mathbf{x}|$，并用下式定义矩阵 $\mathbf{\Gamma}_{\mathbf{x}}(G)$：

\begin{equation*}
\mathbf{T}_{G}\mathbf{x}_{i} = \sum_{j=1}^{d} \mathbf{x}_{j}\mathbf{\Gamma}_{\mathbf{x}}(G)_{ji} \qquad (i = 1 \text{ 到 } d),
\tag{1.3.1}
\end{equation*}

则称 $\mathbf{\Gamma}_{\mathbf{x}}(G)$ 为基 $\langle \mathbf{x}|$ 下表示 $\gamma$ 中代表 $G$ 的矩阵。全部互不相同的矩阵 $\mathbf{\Gamma}_{\mathbf{x}}(G)$ 构成一个矩阵群，它是 $\mathbf{G}$ 在映射 $G \rightarrow \mathbf{\Gamma}_{\mathbf{x}}(G)$ 下的同态像；该同态的核是那些被映到单位矩阵的 $\mathbf{G}$ 的元素。

\textit{例} 1.3.1. 取 $\mathbf{G}$ 为群 $\mathbf{G}_6^2$，其乘法表见表 1.1；该群的几何实现见例 1.2.1。这涉及以 $O$ 为重心的正三角形 $ABC$。令 $\mathbf{OA} = \mathbf{a}$，$\mathbf{OB} = \mathbf{b}$，$\mathbf{OC} = \mathbf{c}$，则 $\mathbf{a} + \mathbf{b} + \mathbf{c} = 0$，且三角形所在平面构成一个二维向量空间 $V$。取其作为表示 $\gamma$ 的基础向量空间，并取 $\mathbf{x}_{1} = \mathbf{b}$，$\mathbf{x}_{2} = \mathbf{c}$ 为基。表示 $\gamma$ 把 $P$、$Q$ 映到 $\mathbf{T}_{P}$、$\mathbf{T}_{Q}$：前者是绕 $O$ 的 120° 逆时针旋转，后者是关于直线 $AO$ 的反射。按定义 $\mathbf{T}_{P}\mathbf{b} = \mathbf{c}$，$\mathbf{T}_{P}\mathbf{c} = \mathbf{a} = -\mathbf{b}-\mathbf{c}$，$\mathbf{T}_{Q}\mathbf{b} = \mathbf{c}$，$\mathbf{T}_{Q}\mathbf{c} = \mathbf{b}$。于是由式 (1.3.1) 得

\begin{equation*}
\mathbf{\Gamma}_{\mathbf{x}}(P) = \begin{pmatrix} 0 & -1 \\ 1 & -1 \end{pmatrix} \qquad
\text{及} \qquad \mathbf{\Gamma}_{\mathbf{x}}(Q) = \begin{pmatrix} 0 & 1 \\ 1 & 0 \end{pmatrix}.
\tag{1.3.2}
\end{equation*}

由于 $\gamma$ 是同态，矩阵群的其余成员可由乘法得出，例如

\begin{equation*}
\mathbf{\Gamma}_{\mathbf{x}}(P^{2}) = \mathbf{\Gamma}_{\mathbf{x}}(P)\mathbf{\Gamma}_{\mathbf{x}}(P) = \begin{pmatrix} -1 & 1 \\ -1 & 0 \end{pmatrix},
\tag{1.3.3}
\end{equation*}

等等。容易验证 $\gamma$ 是忠实表示，但其中含有非幺正矩阵。$\gamma$ 的维数为 2。于是若 $\chi_{\gamma}$ 是 $\gamma$ 的特征标，则由定理 1.3.2 与式 (1.3.2)：$\chi_{\gamma}(E) = 2$，$\chi_{\gamma}(P) = \chi_{\gamma}(P^{2}) = -1$，$\chi_{\gamma}(Q) = \chi_{\gamma}(PQ) = \chi_{\gamma}(P^{2}Q) = 0$。

\textit{例} 1.3.2. 与例 1.2.2 定义的同态 $\theta$ 对应，存在一维矩阵表示 $\mathbf{B}$：$\mathbf{B}(P) = 1$，$\mathbf{B}(Q) = -1$。这是因为 $\mathbf{G}_2^1$ 与由数 1 和 $-1$ 构成的 2 阶乘法群同构。

\textit{例} 1.3.3. 对任何群 $\mathbf{G}$，总存在对称的（或称平庸的）一维表示 $A$：对一切 $G \in \mathbf{G}$，$A(G) = 1$。

由例 1.3.2 与 1.3.3 可见，常常无需指明基础向量空间和基就能找到群的矩阵表示：因为抽象群的同态像可以是矩阵群，它们自然就是矩阵表示。但就我们而言，例 1.3.1 所示的、由定义 1.3.4 产生表示的途径显得更自然，因为我们关心的群都有几何实现，而在这些实现中考察元素对适当向量空间的作用最容易构造表示。

\textbf{定理 1.3.4.} \textit{设 $\langle \mathbf{x}|$ 与 $\langle \mathbf{y}|$ 是 $V$ 的两个基，满足}

\begin{equation*}
\mathbf{y}_{k} = \sum_{i=1}^{d} \mathbf{x}_{i}\mathbf{S}_{ik} \qquad (k = 1 \text{ 到 } d)
\tag{1.3.4}
\end{equation*}

\textit{其中 $\mathbf{S}$ 非奇异，则}

\begin{equation*}
\mathbf{\Gamma}_{\mathbf{y}}(G) = \mathbf{S}^{-1}\mathbf{\Gamma}_{\mathbf{x}}(G)\mathbf{S} \qquad \textit{对一切 } G \in \mathbf{G}.
\tag{1.3.5}
\end{equation*}

也就是说，换基只导致等价的矩阵群 $\mathbf{\Gamma}_{\mathbf{x}}(G)$ 与 $\mathbf{\Gamma}_{\mathbf{y}}(G)$。于是由定理 1.3.1，总可以在 $V$ 中取基 $\langle \mathbf{z}|$，使 $\mathbf{\Gamma}_{\mathbf{z}}(G)$ 成为幺正矩阵群。

\textit{例} 1.3.4. 沿用例 1.3.1 的记号，定义新基 $\langle \mathbf{z}|$：$\mathbf{z}_{1} = \left(\tfrac{1}{3}\sqrt{3}\right)(-\mathbf{b} + \mathbf{c})$，$\mathbf{z}_{2} = -(\mathbf{b} + \mathbf{c})$。由于 $\mathbf{x}_{1} = \mathbf{b}$，$\mathbf{x}_{2} = \mathbf{c}$，得

\begin{equation*}
\langle \mathbf{z}_{1}, \mathbf{z}_{2}| = \langle \mathbf{x}_{1}, \mathbf{x}_{2}| \begin{pmatrix} -\tfrac{1}{3}\sqrt{3} & -1 \\ \tfrac{1}{3}\sqrt{3} & -1 \end{pmatrix}.
\tag{1.3.6}
\end{equation*}

于是由定理 1.3.4，

\begin{equation*}
\mathbf{\Gamma}_{\mathbf{z}}(G) = \mathbf{S}^{-1}\mathbf{\Gamma}_{\mathbf{x}}(G)\mathbf{S},
\tag{1.3.7}
\end{equation*}

其中

\begin{equation*}
\mathbf{S} = \begin{pmatrix} -\tfrac{1}{3}\sqrt{3} & -1 \\ \tfrac{1}{3}\sqrt{3} & -1 \end{pmatrix}.
\tag{1.3.8}
\end{equation*}

由式 (1.3.2)、(1.3.7) 与 (1.3.8) 得

\begin{equation*}
\mathbf{\Gamma}_{\mathbf{z}}(P) = \begin{pmatrix} -\tfrac{1}{2} & -\tfrac{1}{2}\sqrt{3} \\ \tfrac{1}{2}\sqrt{3} & -\tfrac{1}{2} \end{pmatrix} \quad \text{及} \quad \mathbf{\Gamma}_{\mathbf{z}}(Q) = \begin{pmatrix} -1 & 0 \\ 0 & 1 \end{pmatrix},
\tag{1.3.9}
\end{equation*}

即在新基 $\langle \mathbf{z}|$ 下矩阵表示 $\mathbf{\Gamma}_{\mathbf{z}}(G)$ 是幺正的。几何上的原因是：$\mathbf{z}_{1}$ 与 $\mathbf{z}_{2}$ 正交且等长，而算符 $\mathbf{T}_{P}$、$\mathbf{T}_{Q}$ 幺正，因此保持角度与长度。注意按定理 1.3.3，$\mathbf{\Gamma}_{\mathbf{z}}(G)$ 与 $\mathbf{\Gamma}_{\mathbf{x}}(G)$ 有相同的特征标。

\textit{我们今后约定：凡有矩阵表示 $\mathbf{\Gamma}(G)$，向量空间 $V$ 中的基都已取得使矩阵 $\mathbf{\Gamma}(G)$ 为幺正，因此可以使用如下关系：}

\begin{equation*}
\mathbf{\Gamma}(G^{-1})_{ij} = \mathbf{\Gamma}^{*}(G)_{ji}.
\tag{1.3.10}
\end{equation*}

由定理 1.3.1 与 1.3.3，这样做不失一般性——我们只考虑幺正的矩阵表示。

\textbf{定义 1.3.5.} \textit{不变子空间.} 设 $\gamma$ 是 $\mathbf{G}$ 的表示，$\mathbf{T} = \gamma\mathbf{G}$ 是作用在向量空间 $V$ 中的非奇异线性算符群。若

\begin{quote}
(i) $U$ 是 $V$ 的向量子空间，且
(ii) 对一切 $\mathbf{T}_{G} \in \mathbf{T}$ 与一切 $\mathbf{x} \in U$，$\mathbf{T}_{G}\mathbf{x} \in U$，
\end{quote}

则称 $U$ 是 $\mathbf{T}$ 的\textit{不变子空间}（invariant subspace）。

\textbf{定义 1.3.6.} \textit{不可约表示.} 若 $V$ 在 $\mathbf{T}$ 下没有真不变子空间（即除 $V$ 自身与零向量外没有不变子空间），则称 $\gamma$ 为\textit{不可约表示}（irreducible representation）。若在 $\mathbf{T}$ 下存在真不变子空间，则称 $\gamma$ \textit{可约}（reducible）。若 $V$ 能分解为若干子空间的直和，其中每个子空间都在 $\mathbf{T}$ 下不变、并且都是 $\mathbf{G}$ 的某个不可约表示的载体空间，则称 $\gamma$ \textit{完全可约}（completely reducible）。

\textbf{定理 1.3.5.} \textit{有限群 $\mathbf{G}$ 的一切表示都是完全可约的。}

\textbf{定理 1.3.6.} \textit{舒尔引理. 设 $\mathbf{\Gamma}(G)$ 与 $\mathbf{\Gamma}'(G)$ 是 $\mathbf{G}$ 的两个不可约表示，若对一切 $G \in \mathbf{G}$ 有 $\mathbf{\Gamma}(G)\mathbf{S} = \mathbf{S}\mathbf{\Gamma}'(G)$，则要么 $\mathbf{S} = 0$，要么 $\mathbf{S}$ 非奇异且 $\mathbf{\Gamma}(G)$ 与 $\mathbf{\Gamma}'(G)$ 等价。}

以下不可约性判据很有用。

\textbf{定理 1.3.7.} \textit{一个表示不可约，当且仅当与表示的所有矩阵都对易的矩阵只能是单位矩阵的标量倍。}

\textbf{定理 1.3.8.} \textit{设 $\mathbf{G}$ 是阶为 $|\mathbf{G}|$ 的群，元素为 $G_{1}, G_{2}, \ldots, G_{|\mathbf{G}|}$。则 $\mathbf{\Gamma}(G)$ 是不可约表示，当且仅当}

\begin{equation*}
\frac{1}{|\mathbf{G}|}\sum_{i=1}^{|\mathbf{G}|} |\chi_{\gamma}(G_{i})|^{2} = 1
\tag{1.3.11}
\end{equation*}

\textit{其中 $\chi_{\gamma}(G_{i})$ 是 $\mathbf{\Gamma}(G_{i})$ 的特征标。}

\textit{例} 1.3.5. $\mathbf{G}_6^2$ 的表示 $A$、$\mathbf{B}$、$\mathbf{\Gamma}_{\mathbf{z}}(G)$（见例 1.3.2 1.3.4——原书如此）都是不可约的。

由本节的定理和论述应当清楚：求有限群 $\mathbf{G}$ 的全部表示，可以归结为求 $\mathbf{G}$ 的全部互不等价的幺正不可约表示。\textit{从现在起，我们把"幺正不可约表示"简称为"表示"（rep）。}为讨论求有限群 $\mathbf{G}$ 的表示问题，采用如下记号：

\medskip
\noindent\begin{tabular}{|l|p{0.62\textwidth}|}
\hline
$\mathbf{G}$ & 群 \\ \hline
$|\mathbf{G}|$ & $\mathbf{G}$ 的阶 \\ \hline
$G_{j}$ & $j = 1$ 到 $|\mathbf{G}|$，$\mathbf{G}$ 的元素 \\ \hline
$r$ & $\mathbf{G}$ 的类的个数 \\ \hline
$C_{t}$ & $t = 1$ 到 $r$，$\mathbf{G}$ 的类 \\ \hline
$r_{t}$ & $C_{t}$ 中元素的个数 \\ \hline
$\mathbf{\Gamma}^{i}$ & $\mathbf{G}$ 的第 $i$ 个表示的标记 \\ \hline
$d_{i}$ 或 $\dim \mathbf{\Gamma}^{i}$ & $\mathbf{\Gamma}^{i}$ 的维数 \\ \hline
$\mathbf{V}^{i}$ & $\mathbf{\Gamma}^{i}$ 的基础向量空间（载体空间） \\ \hline
\end{tabular}
\medskip

\medskip
\noindent\begin{tabular}{|l|p{0.62\textwidth}|}
\hline
$\langle \mathbf{z}^{i}|$ & $\mathbf{V}^{i}$ 中表示 $\mathbf{\Gamma}^{i}$ 的一组基 \\ \hline
$\mathbf{z}_{p}^{i}$ & $p = 1$ 到 $d_{i}$，基 $\langle \mathbf{z}^{i}|$ 的元素 \\ \hline
$\chi^{i}(G_{j})$ & $G_{j}$ 在 $\mathbf{\Gamma}^{i}$ 中的特征标 \\ \hline
$\chi^{i}(C_{t})$ & 类 $C_{t}$ 中任一元素在 $\mathbf{\Gamma}^{i}$ 中的特征标 \\ \hline
$\mathbf{\Gamma}^{i}(G_{j})_{pq}$ & 表示 $\mathbf{\Gamma}^{i}$ 中 $G_{j}$ 的矩阵表示元的 $(p, q)$ 元 \\ \hline
$\delta_{pq} = \begin{cases}1 & \text{若 } p = q, \\ 0 & \text{若 } p \neq q,\end{cases}$ & \\ \hline
$\delta^{ik} = \begin{cases}1 & \text{若 } \mathbf{\Gamma}^{i} \text{ 与 } \mathbf{\Gamma}^{k} \textit{ 完全相同}, \\ 0 & \text{若 } \mathbf{\Gamma}^{i} \text{ 与 } \mathbf{\Gamma}^{k} \text{ 不等价}.\end{cases}$ & \\ \hline
\end{tabular}
\medskip

\textbf{定义 1.3.7.} \textit{表示的分类}

\begin{quote}
(i) 若 $\Gamma$ 等价于某个实矩阵群，则称表示 $\Gamma$ 属于\textit{第一类}；
(ii) 若 $\Gamma$ 与 $\Gamma^{*}$ 等价但不与任何实矩阵群等价，则称表示 $\Gamma$ 属于\textit{第二类}；
(iii) 若 $\Gamma$ 与 $\Gamma^{*}$ 不等价，则称表示 $\Gamma$ 属于\textit{第三类}。
\end{quote}

\textbf{定理 1.3.9}

\begin{equation*}
\frac{1}{|\mathbf{G}|}\sum_{j=1}^{|\mathbf{G}|} \chi(G_{j}^{2}) = \begin{cases} 1 & \textit{当且仅当 } \Gamma \textit{ 属于第一类，} \\ -1 & \textit{当且仅当 } \Gamma \textit{ 属于第二类，} \\ 0 & \textit{当且仅当 } \Gamma \textit{ 属于第三类。} \end{cases}
\tag{1.3.12}
\end{equation*}

\textit{例} 1.3.6. $\mathbf{G}_6^2$ 的表示 $A$、$\mathbf{B}$、$\mathbf{\Gamma}_{\mathbf{z}}(G)$ 都属于第一类。

\textbf{定理 1.3.10.} \textit{表示的个数 = 类的个数 = $r$。}

\textit{例} 1.3.7. $\mathbf{G}_6^2$ 有三个类（见例 1.2.12），因此三个表示 $A$、$\mathbf{B}$、$\mathbf{\Gamma}_{\mathbf{z}}(G)$ 就是该群仅有的表示——任何其他表示必定与这三者之一等价。

\textbf{定理 1.3.11}

\begin{equation*}
\sum_{i=1}^{r} d_{i}^{2} = |\mathbf{G}|.
\tag{1.3.13}
\end{equation*}

\textit{例} 1.3.8. 作为定理 1.3.11 的例子：$\mathbf{G}_6^2$ 的阶为 6，其表示的维数为 1、1、2。

\textbf{定理 1.3.12.} \textit{矩阵元的正交关系。}

\begin{equation*}
\sum_{j=1}^{|\mathbf{G}|} \mathbf{\Gamma}^{i*}(G_{j})_{pv}\mathbf{\Gamma}^{k}(G_{j})_{qw} = \frac{|\mathbf{G}|}{d_{i}}\delta^{ik}\delta_{pq}\delta_{vw}.
\tag{1.3.14}
\end{equation*}

\textbf{定义 1.3.8.} \textit{特征标表.} 群的\textit{特征标表}（character table）是一个 $(r \times r)$ 方阵，其元为 $\chi^{i}(C_{t})$（$i = 1$ 到 $r$，$t = 1$ 到 $r$）。

注意，由定理 1.3.3 与 1.3.4，特征标表在向量空间 $\mathbf{V}^{i}$ 的换基下不变。但如果给出的不是特征标表，而是完整的矩阵元表 $\mathbf{\Gamma}^{i}(G_{j})_{pq}$——排成 $(|\mathbf{G}| \times |\mathbf{G}|)$ 方阵（$i = 1$ 到 $r$，$p = 1$ 到 $d_{i}$，$q = 1$ 到 $d_{i}$，$j = 1$ 到 $|\mathbf{G}|$）——则必须注明所用的基 $\langle \mathbf{z}^{i}|$，因为该方阵在任何 $d_{i} > 1$ 的 $\mathbf{V}^{i}$ 换基下并非不变。

若把 $(|\mathbf{G}| \times |\mathbf{G}|)$ 方阵视为矩阵（行按 $(ipq)$ 标号、列按 $j$ 标号，记号含义同上），则定理 1.3.12 表明：以 $\sqrt{(d_{i}/|\mathbf{G}|)}\mathbf{\Gamma}^{i}(G_{j})_{pq}$ 为元的矩阵是幺正的。于是立即可得

\begin{equation*}
\sum_{i=1}^{r}\sum_{p=1}^{d_{i}}\sum_{q=1}^{d_{i}} d_{i}\mathbf{\Gamma}^{i}(G_{j})_{pq}\mathbf{\Gamma}^{i*}(G_{m})_{pq} = |\mathbf{G}|\ \delta_{jm}.
\tag{1.3.15}
\end{equation*}

式 (1.3.14) 给出特征标的正交关系：

\textbf{定理 1.3.13.} \textit{特征标的正交关系}

\begin{equation*}
\sum_{i=1}^{r} r_{t}\chi^{i*}(C_{t})\chi^{k}(C_{t}) = |\mathbf{G}|\ \delta^{ik}.
\tag{1.3.16}
\end{equation*}

进而，若把 $(r \times r)$ 方阵视为矩阵（行按 $i$、列按 $t$ 标号），则定理 1.3.13 表明：以 $\sqrt{(r_{t}/|\mathbf{G}|)}\chi^{i}(C_{t})$ 为元的矩阵是幺正的。于是立即可得

\begin{equation*}
\sum_{i=1}^{r} \chi^{i*}(C_{t})\chi^{i}(C_{s}) = \frac{|\mathbf{G}|}{r_{t}}\delta_{ts}.
\tag{1.3.17}
\end{equation*}

\textit{例} 1.3.9. 对 $\mathbf{G}_6^2$，取 $\mathbf{\Gamma}^{1} = A$，$\mathbf{\Gamma}^{2} = \mathbf{B}$，$\mathbf{\Gamma}^{3} = \mathbf{\Gamma}_{\mathbf{z}}(G)$；再取 $G_{1} = E$，$G_{2} = P$，$G_{3} = P^{2}$，$G_{4} = Q$，$G_{5} = PQ$，$G_{6} = P^{2}Q$。则 $C_{1} = \{G_{1}\}$，$C_{2} = \{G_{2}, G_{3}\}$，$C_{3} = \{G_{4}, G_{5}, G_{6}\}$，于是 $r_{1} = 1$，$r_{2} = 2$，$r_{3} = 3$；且 $d_{1} = 1$，$d_{2} = 1$，$d_{3} = 2$。以下对 $\mathbf{\Gamma}^{3}$ 采用例 1.3.4 定义的基 $\mathbf{V}^{3}$ 中由式 (1.3.9) 给出的矩阵。复合标记 $(ipq)$ 按字典序排列：$(111)$，$(211)$，$(311)$，$(312)$，$(321)$，$(322)$。于是下述 $(3 \times 3)$ 幺正矩阵（元为 $\sqrt{(r_{t}/|\mathbf{G}|)}\chi^{i}(C_{t})$，行按 $i$、列按 $t$ 标号）演示了式 (1.3.16) 与 (1.3.17) 的特征标正交关系：

\begin{equation*}
\begin{pmatrix} \dfrac{1}{\sqrt{6}} & \dfrac{1}{\sqrt{3}} & \dfrac{1}{\sqrt{2}} \\[2mm] \dfrac{1}{\sqrt{6}} & \dfrac{1}{\sqrt{3}} & -\dfrac{1}{\sqrt{2}} \\[2mm] \dfrac{2}{\sqrt{6}} & -\dfrac{1}{\sqrt{3}} & 0 \end{pmatrix},
\end{equation*}

而下面的 $(6 \times 6)$ 幺正矩阵（元为 $\sqrt{(d_{i}/|\mathbf{G}|)}\mathbf{\Gamma}^{i}(G_{j})_{pq}$，行按 $(ipq)$、列按 $j$ 标号）演示了式 (1.3.14) 与 (1.3.15) 的矩阵元正交关系：

\begin{equation*}
\left( \begin{array}{cccccc}
\dfrac{1}{\sqrt{6}} & \dfrac{1}{\sqrt{6}} & \dfrac{1}{\sqrt{6}} & \dfrac{1}{\sqrt{6}} & \dfrac{1}{\sqrt{6}} & \dfrac{1}{\sqrt{6}} \\[2mm]
\dfrac{1}{\sqrt{6}} & \dfrac{1}{\sqrt{6}} & \dfrac{1}{\sqrt{6}} & -\dfrac{1}{\sqrt{6}} & -\dfrac{1}{\sqrt{6}} & -\dfrac{1}{\sqrt{6}} \\[2mm]
\dfrac{1}{\sqrt{3}} & -\dfrac{1}{2\sqrt{3}} & -\dfrac{1}{2\sqrt{3}} & \dfrac{1}{\sqrt{3}} & \dfrac{1}{2\sqrt{3}} & \dfrac{1}{2\sqrt{3}} \\[2mm]
0 & -\dfrac{1}{2} & \dfrac{1}{2} & 0 & -\dfrac{1}{2} & \dfrac{1}{2} \\[2mm]
0 & \dfrac{1}{2} & -\dfrac{1}{2} & 0 & -\dfrac{1}{2} & \dfrac{1}{2} \\[2mm]
\dfrac{1}{\sqrt{3}} & -\dfrac{1}{2\sqrt{3}} & -\dfrac{1}{2\sqrt{3}} & \dfrac{1}{\sqrt{3}} & -\dfrac{1}{2\sqrt{3}} & -\dfrac{1}{2\sqrt{3}}
\end{array} \right)
\end{equation*}

\textbf{定理 1.3.14.} \textit{设 $\Gamma$ 是 $\mathbf{G}$ 的任意一个矩阵表示，特征标为 $\chi$。当 $\Gamma$ 经适当的等价变换完全约化成块对角形后，它成为若干表示的直和 $\sum_{i=1}^{r} c_{i}\mathbf{\Gamma}^{i}$，其中}

\begin{equation*}
c_{i} = \frac{1}{|\mathbf{G}|}\sum_{t=1}^{r} r_{t}\chi(C_{t})\chi^{i*}(C_{t}).
\tag{1.3.18}
\end{equation*}

定理 1.3.11 与 1.3.13 有助于确定特征标表；再加上如下定理，计算特征标的任务就成为完全确定的问题。

\textbf{定理 1.3.15}

\begin{equation*}
r_{t}r_{s}\chi^{i}(C_{t})\chi^{i}(C_{s}) = d_{i}\sum_{w=1}^{r} h_{ts,w}r_{w}\chi^{i}(C_{w}),
\tag{1.3.19}
\end{equation*}

\textit{其中 $h_{ts,w}$ 是由式 (1.2.2) 定义的类乘系数。}

\textit{例} 1.3.10. 记 $\chi^{i}(C_{t}) = x_{t}$，$d_{i} = d$，并把例 1.2.12 中的 $r_{t}$ 与 $h_{ts,w}$ 数值代入，则对 $\mathbf{G}_6^2$ 式 (1.3.19) 成为：$x_{1}^{2} = dx_{1}$，$2x_{1}x_{2} = 2dx_{2}$，$3x_{1}x_{3} = 3dx_{3}$，$4x_{2}^{2} = 2dx_{1} + 2dx_{2}$，$6x_{2}x_{3} = 6dx_{3}$，$9x_{3}^{2} = 3dx_{1} + 6dx_{2}$。由定理 1.3.11 知 $d \leqslant 2$。同时满足定理 1.3.13 的三组解为：$d = 1$，$x_{1} = 1$，$x_{2} = 1$，$x_{3} = 1$；$d = 1$，$x_{1} = 1$，$x_{2} = 1$，$x_{3} = -1$；$d = 2$，$x_{1} = 2$，$x_{2} = -1$，$x_{3} = 0$。它们分别对应表示 $A$、$\mathbf{B}$ 与 $\mathbf{\Gamma}_{\mathbf{z}}(G)$。

\textbf{定理 1.3.16.} \textit{外直积. 设 $\mathbf{G} = \mathbf{H} \otimes \mathbf{K}$；若 $\boldsymbol{\Delta}^{i}$（$i = 1$ 到 $r$）与 $\boldsymbol{\Lambda}^{j}$（$j = 1$ 到 $s$）分别是 $\mathbf{H}$ 与 $\mathbf{K}$ 的表示，则 $\mathbf{G}$ 的表示可如下构造。若 $G = (H, K)$，定义 $\mathbf{\Gamma}^{ij}(G) = \boldsymbol{\Delta}^{i}(H) \otimes \boldsymbol{\Lambda}^{j}(K)$，即克罗内克矩阵积，其矩阵元由}

\begin{equation*}
\mathbf{\Gamma}^{ij}(G)_{pq;rs} = \boldsymbol{\Delta}^{i}(H)_{pr}\boldsymbol{\Lambda}^{j}(K)_{qs}
\tag{1.3.20}
\end{equation*}

\textit{给出。于是 $\mathbf{\Gamma}^{ij}(G)$ 是 $\mathbf{G}$ 的表示，让 $i$ 与 $j$ 取遍所有配对即可得到 $\mathbf{G}$ 的全部 $rs$ 个表示。$\mathbf{\Gamma}^{ij}$ 的基是笛卡尔积空间 $\mathbf{V}^{ij} = \mathbf{U}^{i} \otimes \mathbf{W}^{j}$ 中的 $|\mathbf{z}^{i}\mathbf{t}^{j}\rangle$，其中 $|\mathbf{z}^{i}\rangle$ 是 $\mathbf{U}^{i}$ 中 $\boldsymbol{\Delta}^{i}$ 的基，$|\mathbf{t}^{j}\rangle$ 是 $\mathbf{W}^{j}$ 中 $\boldsymbol{\Lambda}^{j}$ 的基。}

由式 (1.3.20) 还可得到 $\mathbf{\Gamma}^{ij}$ 的特征标：

\begin{equation*}
\chi^{ij}(G) = \chi^{i}(H)\chi^{j}(K)
\tag{1.3.21}
\end{equation*}

特别地

\begin{equation*}
\dim \mathbf{\Gamma}^{ij} = (\dim \boldsymbol{\Delta}^{i})(\dim \boldsymbol{\Lambda}^{j}).
\tag{1.3.22}
\end{equation*}

$\mathbf{H} \boxtimes \mathbf{H}$（$\mathbf{H}$ 与自身的内直积）是 $\mathbf{H} \otimes \mathbf{H}$ 中形如 $(H, \dot{H})$ 的元素组成的子群。用定理 1.3.16 的记号，该子群中元素的表示矩阵为

\begin{equation*}
\mathbf{\Gamma}^{ij}(H, \dot{H}) = \boldsymbol{\Delta}^{i}(H) \otimes \boldsymbol{\Delta}^{j}(\dot{H}).
\tag{1.3.23}
\end{equation*}

但在内直积上 $\mathbf{\Gamma}^{ij}$ 是可约的，且等于直和 $\sum_{k=1}^{r} c_{ij,k}\boldsymbol{\Delta}^{k}$（因为 $\mathbf{H} \boxtimes \mathbf{H}$ 与 $\mathbf{H}$ 同构，其表示就是 $\mathbf{H}$ 的表示）。再由定理 1.3.14 与式 (1.3.21) 得

\begin{equation*}
c_{ij,k} = \frac{1}{|\mathbf{H}|}\sum_{p=1}^{|\mathbf{H}|} \chi^{i}(H_{p})\chi^{j}(H_{p})\chi^{k*}(H_{p}),
\tag{1.3.24}
\end{equation*}

其中 $H_{p}$（$p = 1$ 到 $|\mathbf{H}|$）是 $\mathbf{H}$ 的元素。

\textbf{定义 1.3.9.} \textit{克莱布希--戈登分解.} 分解

\begin{equation*}
(\boldsymbol{\Delta}^{i} \boxtimes \boldsymbol{\Delta}^{j})(H) = \boldsymbol{\Delta}^{i}(H) \otimes \boldsymbol{\Delta}^{j}(H) \equiv \sum_{k=1}^{r} c_{ij,k}\boldsymbol{\Delta}^{k}(H),
\tag{1.3.25}
\end{equation*}

其中符号 $\equiv$ 表示定义 1.3.2 意义下的"等价"，而 $c_{ij,k}$ 由式 (1.3.24) 给出，称为表示 $\boldsymbol{\Delta}^{i}$ 与 $\boldsymbol{\Delta}^{j}$ 的内直积 $\boldsymbol{\Delta}^{i} \boxtimes \boldsymbol{\Delta}^{j}$ 的\textit{克莱布希--戈登分解}（Clebsch--Gordan decomposition），系数 $c_{ij,k}$ 称为\textit{克莱布希--戈登系数}（Clebsch--Gordan coefficients），有时也称\textit{维格纳系数}（Wigner coefficients）。注意

\begin{equation*}
c_{ij,k} = c_{ji,k}. \tag{1.3.26}
\end{equation*}
